\begin{center}\begin{tikzpicture}[font=\small,>=stealth,line join=round]
\draw[gray!40,rounded corners=3pt] (0,0) rectangle (14.2,8.3);\node at (7.1,7.700000000000001) {恒零函数：每一步既是反号，也同时保持原值};
\begin{scope}[shift={(2.317391304347826,3.8842857142857143)},x=0.8695652173913043cm,y=1.3142857142857143cm]\draw[->,gray] (-1.4,0)--(5.5,0) node[right]{$x$};\draw[->,gray] (0,-1.7)--(0,1.8) node[above]{$y$};\draw[blue,thick,solid,domain=-1:5,samples=180] plot(\x,{0});\fill[orange!85!black] (0,0) circle(2.2pt);\fill[orange!85!black] (2,0) circle(2.2pt);\fill[orange!85!black] (4,0) circle(2.2pt);\end{scope}\node[anchor=west,align=left,text width=5.0cm] at (8.4,4.0) {$f(x)\equiv0$\\[9pt]$-f(x)=f(x)=0$\\[9pt]任意正数都是周期\\[9pt]不存在最小正周期\\[9pt]此时 $|a|$ 也是周期};
\node[text width=12.6cm,align=center] at (7.1,.65) {非恒零函数移动|a|会在某处改变函数值；恒零函数是这里的例外。};
\end{tikzpicture}\end{center}
